\section{Energy reconstruction and resolution}

\subsection{Measurand definition}

Energy resolution is not a single detector-independent number. For the single-stave response chain, two physically distinct targets are available:
\begin{align}
E_{\rm raw} &\equiv \text{unquenched Geant4 energy deposited in the scintillator},\\
E_{\rm vis} &\equiv \text{Birks-visible/quenched energy used to generate scintillation light}.
\end{align}
A full range-stack reconstruction may later target incident charged-particle kinetic energy from amplitudes and stopping depth. These estimands must not share one bare \(E_{\rm dep}\) label.

For a chosen target \(E_T\), define
\[
r=\frac{E_{\rm reco}-E_T}{E_T}.
\]
A publication result should report median bias, \(\sigmaeight\), RMS, tail fraction and uncertainty intervals on held-out populations, as functions of target energy and relevant phase-space variables.

\subsection{Status of the current held-out result}

\begin{publicationhold}
The current 8.9\% held-out \(\sigmaeight\) result is not a publication deposited-energy resolution. The producer reconstructs \texttt{edep\_scint\_MeV}, which is the Birks-visible variable, while the manuscript/issue called it Geant4 deposited energy. The saturation diagnostic was sign-reversed and clipped to zero, executed-script provenance was incomplete, position transfer was not tested, and the two held-out grid points occupy nearly the same target-energy regime. In addition, the underlying optical grid is superseded by issue \#1303. The plot is therefore quarantined pending a regenerated optical campaign and an explicit \(E_{\rm raw}\) versus \(E_{\rm vis}\) target choice.
\end{publicationhold}

\GatedFigure{figures/gated/edep_reconstruction_heldout.png}{0.90}{Quarantined Figure 11 from the historical held-out reconstruction study.}{GATED: wrong/ambiguous energy semantics and superseded optical campaign.}

\subsection{Required publication analysis}

The replacement study should freeze train/calibration and evaluation populations before fitting; include more than one effective target-energy regime; validate particle-species transfer without truth-label tuning; test longitudinal/transverse position and angle support; compare transparent linear/monotonic baselines before any ML model; and propagate optical/SiPM model nuisance families separately from finite-simulation statistics. If a resolution parametrisation is fitted, the available energy grid must identify its terms.

A full-stack incident-energy estimator may use the amplitude vector \(\mathbf A=(A_{B2},A_{B4},A_{B6},A_{B8})\) and stopping depth, but it remains downstream of material-budget, trigger/selection and per-channel response closure.
